% ch.tex

%%%%%%%%%%%%%%%
\begin{frame}{}
  \begin{center}
    \teal{\Large Continuum Hypothesis (CH)}
  \end{center}

  \[
    \red{\fbox{\text{\LARGE $\nexists A: \aleph_{0} < |A| < \mathfrak{c}$}}}
  \]
\end{frame}
%%%%%%%%%%%%%%%

%%%%%%%%%%%%%%%
\begin{frame}{}
  \begin{columns}
    \column{0.50\textwidth}
      \fig{width = 0.90\textwidth}{figs/dangerous-knowledge-bbc-cover}
    \column{0.50\textwidth}
      \fig{width = 0.80\textwidth}{figs/dangerous-knowledge-people}
  \end{columns}

  \begin{center}
    {\leftpointright\; \purple{Dangerous Knowledge (22:20; BBC 2007)}}
  \end{center}

  \pause
  Independence from ZFC: \\[1pt]
  \begin{description}
    \item[Kurt G\"odel (1940)] CH cannot be disproved from ZF.
    \item[Paul Cohen (1964)] CH cannot be proven from the ZFC axioms.
  \end{description}
\end{frame}
%%%%%%%%%%%%%%%

%%%%%%%%%%%%%%%
\begin{frame}{}
  \fig{width = 0.55\textwidth}{figs/cantor-monumento-labelled}

  \begin{quote}
    \begin{center}
      \blue{``das wesen der mathematik liegt in ihrer freiheit''} \\[8pt]

      \red{``The essence of mathematics lies in its freedom''}
    \end{center}
  \end{quote}
\end{frame}
%%%%%%%%%%%%%%%

%%%%%%%%%%%%%%%
\begin{frame}{}
  \begin{quotation}
    ``这对我来说是最值得钦佩的数学理智之花，\\[6pt]
    也是在纯粹理性范畴中人类活动所取得的最高成就之一。''\\[5pt]
    \hfill --- 希尔伯特
  \end{quotation}

  \begin{columns}
    \column{0.33\textwidth}
      \fig{width = 0.70\textwidth}{figs/engines-of-logic}
    \column{0.33\textwidth}
      \fig{width = 0.70\textwidth}{figs/logic-history}
    \column{0.33\textwidth}
      \fig{width = 0.70\textwidth}{figs/set-theory}
  \end{columns}
\end{frame}
%%%%%%%%%%%%%%%